\chapter{Development: Why the Clean Threshold Story Fails, and What Replaces It}
\label{ch:development}

\begin{plainlanguage}
\textbf{In plain terms:} A child's ability to \emph{tell} two animals apart and their ability to \emph{say} the right word for each turn out to run on separate, loosely connected tracks, often years apart — a toddler saying "doggie" for a cow may already perceptually distinguish the two. This forces tracking perception, comprehension, and speech as separate systems, with a strict rule against inventing such splits just to save a failing prediction.
\end{plainlanguage}

\section{The naive prediction, again stated in advance}

\begin{naivemodel}
A second domain, chosen for being about as unlike flood-control infrastructure as a case study can be: how children acquire and differentiate word categories. The naive hypothesis under test here is structurally identical to Chapter~\ref{ch:room-for-river}'s, transplanted: a child maintains a single, unified category structure (a ``distinction graph'') that accumulates anomalies --- instances that don't fit existing categories --- until accumulated anomaly forces a closure-expanding reorganization, splitting an overextended category into two. Applied to the well-documented phenomenon of childhood overextension (a child calling cows, horses, and sheep all ``dog''), the prediction is that new categories (dog vs.\ cow, as separate words) should emerge specifically when anomaly accumulation against the single existing category graph exceeds what the graph can absorb.
\end{naivemodel}

\section{Where it breaks}

This prediction does not survive contact with the developmental-psychology literature, and like the institutional case, it breaks in a way that is more useful than confirmation would have been.

First: production and comprehension dissociate. In preferential-looking studies, children showed no looking-time preference distinguishing ``dog'' from ``cat'' on anomalous trials (where no matching referent was present), yet did show a preference when ``cow'' was the requested label --- and critically, prior overextension \emph{in production} (the child saying ``dog'' for a cow) was not diagnostic of the child's underlying comprehension. A child can be overextending a word in speech while their comprehension already discriminates the categories that speech is failing to mark. This means there is no single ``the child's category graph'' whose saturation level can be tracked --- there are at minimum two systems, production and comprehension, updating on different schedules and potentially containing different information at the same chronological age.

Second: in several documented cases, perceptual category differentiation precedes lexical differentiation rather than being generated by it. Eye-tracking work on infant categorization shows that visual category representations distinguishing cats from dogs, built from head-region cues, are established in early infancy --- well before the one-to-two-and-a-half-year window in which overextension in speech is typically observed. The perceptual distinction (cat $\neq$ dog) is often already in place before the lexical failure (calling both ``doggie'') even begins; the lexical layer is not generating a new distinction under anomaly pressure so much as catching up to one that already existed in a different subsystem.

\section{The repair}

\begin{repairbox}
\begin{definition}[Subsystem-Indexed Repair]
Rather than tracking a single global repair process $R_n(X)$ for ``the child,'' index repair by subsystem: $R_n^{(k)}(X)$ for subsystem $k$ --- for example, perception, lexicon, or production.
\end{definition}

This move is not free, and stating the cost of it plainly matters more here than almost anywhere else in the book, because it is exactly the kind of move that can quietly become unfalsifiable if not constrained. Splitting a global process into per-subsystem processes whenever a prediction fails is a standard and tempting way to rescue a theory from any disconfirming evidence whatsoever --- simply declare a new subsystem boundary wherever the data doesn't fit, and no observation can ever count against the framework.

\begin{principle}[Anti-Vacuity]
A subsystem split into index $k$ is admissible only if there exists an independent measurement method $M_k$ for that subsystem --- established by existing methodology for reasons unrelated to rescuing this particular prediction, not introduced solely because the global version failed.
\end{principle}

This is not a mathematical claim and has no proof; it is a methodological commitment, and it is satisfied here, not assumed: the perception/comprehension/production split in this chapter is not invented to save the anomaly-threshold hypothesis. It is borrowed directly from existing developmental-psychology methodology --- the preferential-looking paradigm exists, and was developed, precisely because researchers already suspected production and comprehension could dissociate, for reasons that have nothing to do with this book's elimination formalism. The split earns its use here because it was independently motivated before this chapter needed it.
\end{repairbox}

\section{A second, smaller dissociation}

The within-subsystem story is also not clean. Even confined to comprehension and production separately, the literature warns against treating either as a single unitary process: the relevant developmental and self-regulatory literatures generally caution against assuming a shared core mechanism across what look, behaviorally, like related capacities --- a theme that recurs with sharper formal content in Part~IV, where the same warning becomes load-bearing for the recursive compression argument rather than a side note.

\section{What survives}

Not the clean anomaly-threshold hypothesis. What survives, and is consistent with --- though again not independently confirmed by --- Chapter~\ref{ch:room-for-river}'s coupling-lag conjecture: subsystems within a single developing child can be loosely coupled enough to update on substantially different timescales, with perception sometimes leading production by a year or more, and apparent ``category collapse'' at one level (the spoken word ``doggie'' applied indiscriminately) can coexist with a fully maintained distinction at another (the eye-tracked perceptual discrimination between the two animals). This recasts what looked like a developmental error --- overextension --- as something closer to economical repair at the lexical/production interface specifically, with no failure at all occurring at the perceptual level.
